\section{Maximum Bracket Size pada Team Event}

Pada \textit{team event}, peserta yang digunakan untuk membentuk bracket eliminasi bukan lagi individu, melainkan \textbf{jumlah tim yang terbentuk}. 
Setiap tim biasanya terdiri dari beberapa atlet (misalnya 3 atlet per tim pada kompetisi panahan standar).

\subsection{Pembentukan Tim}

Jika terdapat $N$ atlet dan setiap tim terdiri dari $k$ atlet, maka jumlah maksimum tim yang dapat terbentuk adalah:

\[
T = \left\lfloor \frac{N}{k} \right\rfloor
\]

dengan:

\begin{itemize}
\item $N$ = jumlah atlet
\item $k$ = jumlah atlet per tim
\item $T$ = jumlah tim maksimum
\item $\lfloor \cdot \rfloor$ = fungsi pembulatan ke bawah
\end{itemize}

\textbf{Contoh:}

Jika terdapat $N = 60$ atlet dan setiap tim terdiri dari $k = 3$ atlet:

\[
T = \left\lfloor \frac{60}{3} \right\rfloor = 20
\]

Secara teoritis dapat terbentuk maksimal \textbf{20 tim}.  
Namun dalam praktiknya jumlah tim bisa lebih sedikit karena pembentukan tim biasanya berdasarkan klub atau negara.

\subsection{Penentuan Ukuran Bracket Tim}

Ukuran bracket eliminasi untuk tim juga mengikuti bilangan pangkat dua:

\[
2, 4, 8, 16, 32
\]

Ukuran bracket ditentukan dengan rumus:

\[
\text{Bracket Size} = 2^{\lceil \log_2(T) \rceil}
\]

dengan:

\begin{itemize}
\item $T$ = jumlah tim yang berpartisipasi
\item $\lceil \cdot \rceil$ = fungsi pembulatan ke atas
\end{itemize}

\subsection{Contoh}

Misalkan hanya terbentuk \textbf{12 tim} dari hasil qualification.

\[
\log_2(12) \approx 3.58
\]

\[
\lceil 3.58 \rceil = 4
\]

Sehingga ukuran bracket yang digunakan adalah:

\[
2^4 = 16
\]

Artinya bracket eliminasi memiliki \textbf{16 slot tim}.

\subsection{Perhitungan Bye}

Jika jumlah tim lebih kecil dari ukuran bracket, maka slot kosong akan menjadi \textbf{bye}.

\[
\text{Bye} = \text{Bracket Size} - T
\]

Contoh:

\[
16 - 12 = 4
\]

Sehingga terdapat \textbf{4 bye} pada ronde pertama.

Bye biasanya diberikan kepada tim dengan \textbf{ranking qualification tertinggi}.

\subsection{Ringkasan}

\begin{itemize}
\item Team event menggunakan jumlah \textbf{tim} sebagai dasar pembentukan bracket.
\item Ukuran bracket tetap mengikuti bilangan pangkat dua.
\item Maximum bracket size ditentukan dengan fungsi $2^{\lceil \log_2(T) \rceil}$.
\item Jika jumlah tim lebih kecil dari ukuran bracket, maka slot kosong menjadi \textbf{bye}.
\end{itemize}